\problemname{Drawing Borders}

Somewhere in the great North American plains live the tribes of
chiefs \emph{Blue Eagle}, \emph{Red Beaver}, and \emph{Green Serpent}.
Their population is scattered over numerous villages all over the
land and conflict arises whenever members of different tribes meet
while traveling across the plains.

To put an end to the constant animosities the chiefs have decided that
the land should be divided between the tribes so that they can avoid
each other when moving between villages belonging to the same tribe.
More precisely, they want to construct two border fences -- thus dividing
the land into three regions -- such that two villages lie in the same
region precisely when they belong to the same tribe.

The villages are represented by points in the Euclidean plane that are
colored blue, red or green, depending on the tribe, and the fences
should be drawn in the form of two polygons. The polygons may not
touch or intersect themselves or each other and none of the points
may lie on their boundary. (Make sure to read the constraints in the
Output section!)

\begin{figure}[!h]
\centering
\includegraphics[width=8cm]{../data/sample/1.png}
\caption{Illustration of the sample.}
\end{figure}

\section*{Input}

The input consists of:

\begin{itemize}
\item one line with an integer $n$ ($3 \le n \le 100$), the number of
villages.
\item $n$ lines, each with three integers $x, y, c$ ($-1\,000 \le x, y \le
1\,000, 1 \le c \le 3$), representing a village at coordinates
$(x,y)$ of color $c$ ($1=$ blue, $2=$ red, $3=$ green). All positions
are unique. There is at least one village of each color.
\end{itemize}

\section*{Output}

If there is no solution, print \texttt{impossible}. Otherwise, print
the two polygons in the following format:

\begin{itemize}
\item one line with an integer $m$ ($3 \le m \le 1\,000$), the number
of vertices of the polygon.
\item $m$ lines, each with two real numbers $x, y$ ($-3\,000 \le x, y \le
3\,000$), the vertices of the polygon in either clockwise or
counter-clockwise order. The numbers may be given with up to five
decimal places (additional places will be rounded off).
\end{itemize}

